同一个矩阵反复作用，是数据工作中出现频率最高的动作之一：Markov 链一步步转移，PageRank 一轮轮分发权重，循环网络一层层传递隐状态，梯度下降一次次施加同一个迭代算子。这些问题都不问``乘一次会怎样''，而问``乘一百次以后是什么样''。直接把矩阵乘上去，坐标每一步都在互相混合，很快就既看不出结构也控制不住数量级。

本章的全部内容来自一个换角度的念头：在 \(A\) 的作用下，多数方向都会被拧过去，但也许存在一些方向拧不动——作用之后仍落在原来那条直线上，只是长度变了。沿这样的方向，矩阵退化成一个数，重复作用就是这个数反复自乘。

\begin{center}
\begin{tikzpicture}[x=1cm,y=1cm,>=stealth,font=\small]
  \draw[black!45,fill=blue!7,rounded corners=2pt] (0,0) rectangle (3.6,1.5);
  \node[align=center] at (1.8,0.75) {标准坐标\\\(x_{k+1}=Ax_k\) 分量耦合};
  \draw[black!45,fill=blue!16,rounded corners=2pt] (5.4,0) rectangle (9.6,1.5);
  \node[align=center] at (7.5,0.75) {特征坐标\\\(n\) 个独立标量 \(\lambda_i^{\,k}\)};
  \draw[->,thick] (3.75,1.05) -- (5.25,1.05);
  \node[above,font=\footnotesize] at (4.5,1.1) {\(S^{-1}\)};
  \draw[<-,thick] (3.75,0.45) -- (5.25,0.45);
  \node[below,font=\footnotesize] at (4.5,0.4) {\(S\)};
  \node[align=center,black!70] at (7.5,-0.85) {长期行为由最大模的 \(\lambda\) 主导};
\end{tikzpicture}

\emph{整章的骨架：换一组坐标，把纠缠的迭代拆成互不干扰的标量幂；难点全部集中在这组坐标是否存在。}
\end{center}

以 \(A=\begin{pmatrix}2&1\\1&2\end{pmatrix}\) 为例。标准坐标里 \((x,y)\) 变成 \((2x+y,\,x+2y)\)，两个分量牵扯在一起；但沿 \((1,1)^{\trans}\) 的向量被整齐放大三倍，沿 \((1,-1)^{\trans}\) 的向量纹丝不动。把这两条斜线当坐标轴，矩阵就成了 \(\diag(3,1)\)。矩阵没变，变的是记录它的坐标。

四篇沿着这条主线推进。第一篇把``不转向的方向''变成可以求解的对象，并说明特征多项式虽然定义了特征值，却不是数值上求特征值的办法。第二篇假定这样的方向恰好够用，于是任意初始向量都能拆成独立演化的分量，长期行为由最大模的特征值说了算——Markov 链的平稳分布与谱间隙、PageRank、幂迭代、循环网络的梯度消失与爆炸，都是这一条的不同外衣。第三篇拿掉这个前提：方向不够时会出现多项式增长，临界情形因此没有安全余量。第四篇把矩阵幂与矩阵指数收进同一个框架，连续时间的线性系统与控制由此接上。

\begin{itemize}
\tightlist
\item \hyperref[sec:03-Linear-Algebra/05-Eigenvalues-and-Eigenvectors/01]{``不变方向与特征值：寻找矩阵最自然的坐标轴''}——把``哪些方向不转''写成可解的方程，并区分几何重数与代数重数；
\item \hyperref[sec:03-Linear-Algebra/05-Eigenvalues-and-Eigenvectors/02]{``对角化与动力系统：把耦合演化变成独立缩放''}——判断长期行为不必迭代，看谱半径与谱间隙就够；
\item \hyperref[sec:03-Linear-Algebra/05-Eigenvalues-and-Eigenvectors/03]{``对角化失败与 Jordan 形：独立特征方向不够时怎么办''}——特征值全在单位圆上也可能发散，缺陷结构给出多项式因子；
\item \hyperref[sec:03-Linear-Algebra/05-Eigenvalues-and-Eigenvectors/04]{``矩阵函数：从多项式到指数''}——\(f(A)\) 只由 \(f\) 在谱上的取值（必要时加导数）决定。
\end{itemize}

首次阅读把重心放在前两篇：不变方向是什么，以及一组完整的不变方向怎样把迭代拆开。第三篇的 Jordan 基手算可以跳过，但``几何重数不足会产生多项式增长''这个结论要带走。第四篇适合当作回读入口，写线性微分方程或做连续时间建模时再翻。若已经在别处见过谱定理，请注意本章讨论的是一般方阵，特征向量并不天然正交；正交的保证要到\hyperref[ch:021]{``对称矩阵、二次型与正定性''}才拿得到。想知道这些量在浮点数下怎么算出来，入口是\hyperref[sec:08-Numerical-Analysis/05-Numerical-Linear-Algebra/05]{``从幂迭代到 QR 特征值算法''}。
